\documentclass[12pt]{scrartcl}
\usepackage{graphicx}
\usepackage[default]{opensans}
\usepackage{sfmath} % sans font also for math
 \usepackage[binary-units = true]{siunitx}

% defining the paper layout that no text overlaps with the header
\usepackage[
  top=35mm,
  headheight=25mm,
  headsep=3mm,
  bottom=30mm,
  left=25mm,
  right=25mm
]{geometry}

\usepackage{latexsym}
\usepackage[centertags]{amsmath}
\usepackage{amssymb}

% custom header and footpage
\usepackage{scrpage2}
\pagestyle{scrheadings} % you have to set the custom layout
\ihead{D4.2: Long Pulse Laser -- Matter Interaction} % left head
\ohead{\includegraphics[height=25mm]{EUCALL.png}}
\ifoot{\includegraphics[height=13.4mm]{EU.png}} % left foot
\cfoot{
  \begin{minipage}{100mm}
    \begin{scriptsize}
      \normalfont{This project has received funding from the}
      \textit{European Union’s Horizon 2020 research and innovation programme}
      \normalfont{under grant agreement No 654220.}
    \end{scriptsize}
  \end{minipage}
} % center foot
\ofoot{} % right foot
\chead{}

\usepackage{booktabs} % more eye candy for tables

% sophisticated linking of references in the pdf and setting some options
\usepackage{url}                                                  % for correct typesettings of URLs
\usepackage{hyperref}                                             % for sophisticated linking of urls, dois, pictures, tables, etc.
\hypersetup{
    unicode=true,                                                 % non-Latin characters in Acrobat’s bookmarks
    pdftoolbar=true,                                              % show Acrobat’s toolbar?
    pdfmenubar=true,                                              % show Acrobat’s menu?
    pdffitwindow=false,                                           % window fit to page when opened
    pdfstartview={FitH},                                          % fits the width of the page to the window
    pdftitle={D4.1: Short pulse laser-matter interaction simulations},                                        % title
    pdfauthor={C. Fortmann-Grote},                                           % author
    pdfsubject={EUCALL WP4 (SIMEX) deliverable 4.1},                             % subject of the document
    pdfcreator={pdflatex},                                         % creator of the document
    pdfkeywords={EUCALL, SIMEX, simulations, PIC, XFEL, scattering},                                         % list of keywords
    pdfnewwindow=true,                                            % links in new PDF window
    colorlinks=true,                                              % false: boxed links; true: colored links
    linkcolor=blue,                                                % color of internal links (change box color with linkbordercolor)
    citecolor=blue,                                                % color of links to bibliography
    filecolor=blue,                                               % color of file links
    urlcolor=blue                                                 % color of external links
}


% Zeilenabstand
\renewcommand{\baselinestretch}{1.2}


%\renewcommand*\chapterpagestyle{scrheadings} % otherwise chapter start is plain ;-); works only for KOMA script; plain LaTeX works different

%%%%%%%%%%%%%%%%%%%%%%%%%%%%%%%%%%%%%%%%%%%%%%%
%   BIBLIOGRAPHY SETTINGS                     %
%%%%%%%%%%%%%%%%%%%%%%%%%%%%%%%%%%%%%%%%%%%%%%%
\usepackage[bibstyle=nature,sorting=none,=maxnames=1000,eprint=false,
defernumbers=true, backend=biber]{biblatex}
\usepackage{hyperref}

\renewcommand*\finalnamedelim{, and\addspace}
\DeclareNameAlias{sortname}{last-first}
\renewcommand{\newunitpunct}{, }

\AtEveryBibitem{%
  \clearfield{day}%
  \clearfield{month}%
  \clearfield{endday}%
  \clearfield{endmonth}%
  \clearfield{issn}%
  \clearfield{issue}%
}
%convert titles to hyperlinks using doi
\ExecuteBibliographyOptions{doi=false} \newbibmacro{string+doi}[1]{%
  \iffieldundef{doi}{#1}{\href{http://dx.doi.org/\thefield{doi}}{#1}}}
  \DeclareFieldFormat*{title}{\usebibmacro{string+doi}{\mkbibemph{#1}}}

%\addbibresource{/home/grotec/Documents/LiteratureDB/bibtex/jabref.bib}
\addbibresource{references.bib}
%%%%%%%%%%%%%%%%%%%%%%%%%%%%%%%%%%%%%%%%%%%%%%%
% END BIBLIOGRAPHY SETTINGS                   %
%%%%%%%%%%%%%%%%%%%%%%%%%%%%%%%%%%%%%%%%%%%%%%%

\begin{document}
\makeatletter
\begin{titlepage}
\thispagestyle{scrheadings}
\begin{center}
$~$\\
\vspace{2cm}
\Huge{\textbf{WP 4 -- SIMEX\\[1cm]
Design Report and Advanced Simulation Software for Long Pulse Laser -- Matter Interaction}}\\[5mm]
\vspace{2cm}
\large{
Carsten Fortmann-Grote, Alexander Andreev, Richard Briggs,\\ Michael Bussmann,
  Axel Huebl, Thomas Kluge,\\
 S. Pascarelli, Ashutosh Sharma, and Adrian P. Mancuso\\
 }
\vspace{1cm}
\@date
\end{center}
\vfill%
\includegraphics[width=\textwidth]{./PartnerLogos.pdf}
\normalfont
\end{titlepage}
\makeatother

\tableofcontents

\begin{abstract}
  \noindent%
  \textbf{Abstract} -- We present a design for integrated simulations of an x-ray absorption
  experiment probing dynamically compressed states of matter at tens of GPa pressure generated by
  high energy ($\approx 30\,\text{J}$) laser pulses interacting with solid matter.
  We model the laser-matter interaction using 1-D and 2-D radiation-hydrodynamic simulations to predict
  the thermodynamical state of the compressed matter over the course of the x-ray pulse. These simulations
  are embedded in the software framework \texttt{simex\_platform} to allow seamless integration into
  start-to-end simulations of an entire beamline experiment including x-ray pulse propagation in the beamline and
  detector response. The involved simulation codes are briefly presented.
\end{abstract}
%
\section{Introduction}
Long pulse laser systems are capable of reaching intensities of up to $\sim 10^{12}$ -- 10$^{15} \text{ W}/\text{cm}^2$ with pulse lengths of
greater than a few nanoseconds. The pulse length and focal spot of the laser can depend heavily on the sample material and target package.
Generally, focal spot sizes range between $\sim$ 100 $\mu$m up to 1 mm. With these specifications, the laser energies now required to reach the
highest intensities are on the order of hundreds of joules.

Temporally shaped laser pulses interacting with overdense target material create a rapidly expanding plasma that can generate a shockwave in
the ablating material. This shock wave can travel through the target at several km/s and compress the sample material to pressures reaching
several hundred GPa. (For comparison, the Earth's core pressure is $\sim$ 330 GPa). By using a ramp temporal pulse, where the laser intensity
is slowly increased, the temperatures generated remain much cooler than during rapid shock compression and the solid state of matter can be
studied up to several TPa (1 TPa = 1000 GPa). Such states of matter are often referred to as Warm Dense Matter (WDM).

High-energy laser facilities are now becoming increasingly in demand at new generation x-ray light sources (X-FELs, 3rd generation synchrotron
sources). The potential of long pulse laser experiments are already yielding extreme conditions of matter well beyond the reach of static
high-pressure techniques and combinations with X-ray techniques provide enticing experiments at new extreme states of matter.

We design here simulations for a prototypical experiment combining a high-energy laser system and x-ray radiation from a 3rd generation
synchrotron. The laser shock-compresses a tailored solid density target and the compressed matter is subsequently probed by x-rays. X-ray
absorption spectroscopy (XAS) allows to monitor the compression and to characterize the electronic and structural states in the target. Through
variation of the delay between optical laser pulse and x-ray probe pulse, time resolved data is obtained.

\begin{table}[h]
  \centering
  \begin{tabular}{l|l}
    \hline
    \hline
    Photon energy & 5 keV - 28 keV \\
    Polarization & horizontal \\
    Focus size (fwhm) & 3 $\mu$m (5 keV) - 50 $\mu$m (28 keV)\\
    Pulse length & 100 ps\\
    Flux  & $1\times 10^{14}$ photons/s (7 keV) - $4\times 10^{13}$ photons/s (28 keV) \\
    \hline
    \hline
  \end{tabular}
  \caption{Principal x-ray parameters for shock compression studies at ESRF
  (beamline ID24)}
  \label{tab:esrf_parameters}
\end{table}
\section{Photon-Matter interaction}
\subsection{Optical Laser}
Nearly all of the physical processes of long pulse laser-matter interactions are described by partial differential
equations that can be solved with careful construction of numerical simulation codes. Here, the laser-matter
interactions are modelled with the 1D radiation-hydrodynamics computer code ``Esther'' \cite{Colombier2005}
or the 2D code ``MULTI2D'' \cite{Ramis2009}.

Table~\ref{tab:esrf_long_pulse} summarizes the main optical pump laser
parameters for shock compression experiments at the ESRF.

\begin{table}[h]
  \centering
  \begin{tabular}{l|l}
    \hline
    \hline
Wavelength & 1053 nm \\
Pulse duration & 4 ns - 15 ns \\
Intensity & $2\times 10^{13}$ W/cm$^2$ \\
Pulse energy & 30 J \\
Focus size (fhwm) & 100 $\mu\text{m}^{2}$ - 300 $\mu\text{m}^{2}$  \\
Temporal profile  & flat top \\
Spatial profile & Gaussian in x,y \\
    \hline
    \hline
  \end{tabular}
  \caption{ESRF optical laser parameters for long pulse laser-matter
  simulations.}
  \label{tab:esrf_long_pulse}
\end{table}

A call for tender will be issued in 2016 for a new high-energy laser system to be installed at the ESRF, with increased energy (200 J) and
frequency doubling crystals.

For the 1D hydrocodes, a moving mesh is applied in an arbitrary Lagrangian-Eulerian framework where the coordinates of
the mesh contain the variables from which density can be calculated (mass of
each zone is fixed). The most important variables
considered in these codes are those associated with the extreme conditions generated by the high-power lasers: pressure (density),
temperature and velocity. Feedback of these hydrocode outputs are crucial in the design and implementation of laser shock/ramp
experiments with X-ray interactions.

The X-ray probe duration, during laser shock compression, can range from nanosecond down to femtosecond exposures.
X-ray pulses for shock compression studies at the ESRF are $\approx$ 100 ps
long. Table~\ref{tab:esrf_parameters} lists the most important optical laser
parameters.

It is crucial to have an understanding of shock transit times so that accurate timing of the X-ray probe,
with respect to the laser initialisation, can be made. The final density state reached by a shock, as calculated
from the hydrocode simulation packages, can define the shock velocity.
\subsubsection{Esther}
The Esther hydrocode was written by Patrick Combis and Laurent Videau of the CEA, Paris, France.
\footnote{The Esther hydrocode is currently presented in French only.}
A license to use Esther can be obtained by requesting access (via email) from the authors (for academic use only).

As a necessary step for integration of hydrocode modelling in SIMEX, we have
developed interfaces Esther. An interactive input file generator sets up
the configuration files for the hydrocode based on user provided information
about the target and the laser pulse. After the simulation completed, a
converter reformats the hydrocode data (pressure, density, temperature, and
shock velocity as function of space and time) from the native format into an
openPMD conform hdf5 file. In this context, it is worth mentioning, that a new
openPMD standard extension is currently being developed to enable facilitate the
usage of openPMD as a standard format for rad-hydro codes.

\subsubsection{MULTI}
In the MULTI2D hydrocode \cite{Ramis2009},
the hydrodynamic equations solved by the code are combined with a multigroup method for radiation transport.

\subsection{X-ray matter interaction}
In the proposed experimental setup, discussed in detail further below, the shock compressed matter is diagnosed by
X-ray Absorption Fine Structure (XAFS) spectroscopy and radiography. Both techniques can be modeled with appropriate simulation tools as
described in the following.

In some cases the particle velocity of the interface between sample and window cannot be obtained directly as strong shocks can transform the transparent window to an opaque
material. In those instances only the free surface velocity can be obtained by the velocity diagnostics and an iterative Lagrangian analysis (ILA) must be used to determine the sound
speed; stress-density information can then be calculated \cite{Rothman2006}.  Using the ILA, the velocity-time history during the compression can be obtained from the free surface measurements and from knowledge of the optical laser pulse to x-ray pulse delay, the exact state probed by the x-ray is obtained.
\subsection{XAFS}
The XAFS signal reflects the near order of the shock compressed matter. The absorption edge is modulated by the ion-ion correlation function
$\chi(\vec q)$.

\section{Experimental scenario}
\begin{figure}[ht]
  \begin{center}
    \includegraphics[width=1.0\textwidth,angle=0,clip]{torchio2016_setup-crop.pdf}
  \end{center}
  \caption{Schematic rendering of a laser driven shock compression experiment
  and x-ray absorption spectroscopy (a). XANES signal as function of time
indicating structural phase transition from bcc to hcp phase (b). 1-D
Radiation-hydrodynamics simulation of the shock
propagation and pressure dynamics during laser-matter interaction (c).
Figure taken from Ref.~\cite{Torchio2016}.}
  \label{fig:exp_setup}
\end{figure}
%
Fig.~\ref{fig:exp_setup} shows the experimental setup, optical laser and x-ray
probe pulse timings (a), measured XANES spectra (b) and 1-D radiation-hydrodynamic
simulations of shock propagation through the a sandwiched iron target.
The figure is taken from Ref.~\cite{Torchio2016}. The drive pulse compresses the
target as modeled by radiation hydrodynamics simulations reaching pressures in
the range of 40 -- 60 GPa. The x-ray pulse of $\approx 100\,\text{ps}$ probes
the compressed target after a variable time delay. Transmitted x-ray photons are
detected in a position sensitive Ge detector. The position of the detected
photon can be converted into photon energy employing the dispersion curve of the
focussing mirror.

\marginpar{Should discuss novel aspects in a follow-up experiment}

\section{Application: Pulse optimization via feedback loop}
Radiation-hydrodynamic simulations are an indispensable tool to estimate the
achievable level of compression and pressures in a dynamic compression
experiment. Tailored (ramp) pulses can be employed to generate low-isentrope
compressions, i.e. reaching high densities at low temperatures, below the
principal Hugoniot and below the melting temparature. This technique makes
accessible the range of strongly compressed solid matter and study pressure
induced structural and electronic phase transitions.

To facilitate pre-shot simulations and pulse optimization, we will utilize the
simulation capabilities of \textit{simex\_platform} inside an optimization loop (Fig.~\ref{fig:lp_optimization_schem}).
Using an initial guess for the (temporal) pulse shape, we run the
radiation-hydrodynamics simulations and monitor the evolution of density,
pressure, and temperature over the pulse duration. This information is then fed
to an algorithm that sets up a new, improved pulse shape based on given criteria
like maximum allowable temperature or minimal compression ratio, to run the next
iteration of the process.
The pulse shaping tool can be implemented in various
ways including nonlinear multidimensional optimization of pulse parameters or genetic algorithms.

\begin{figure}[ht]
  \begin{center}
    \includegraphics[width=1.0\textwidth,angle=0,clip]{lp_optimization_scheme-crop}
  \end{center}
  \caption{Optimization of pulse shapes for low isentrope compression. The
  output of radiation-hydrodynamics simulations is fed into a pulse shape
generator to optimise the laser-matter interaction for maximum compression at
low temperature increase. Once an optimized pulse is generated, the simulation
continues with calculating the x-ray -- matter interaction and synthesizing the
diagnostic signals}
  \label{fig:lp_optimization_schem}
\end{figure}

\section{Synergy}

Although our pulse optimization scheme is  independent of the technological and
physical realization of pulse generation and control, it is desirable to
incorporate the actual pulse generation mechanism, it's capabilities as well as
it's limitations into the simulation workflow.  This opens an opportunity for
a cross-workpackage collaboration with the EUCALL workpackage 7 (``Pulse Control and
ChAracterization -- PUCCA'')).

The development of an open data standard for radiation-hydrodynamics simulation is of high value for the scientific community
also beyond EUCALL. It will facilitate integration of radiation-hydrodynamics codes into simulation frameworks such as
\texttt{simex\_platform}, benchmarking of various codes against each other, against other methods (e.g. particle-in-cell simulations), and
against experimental data.

\section{Summary}
We outlined the design for simulations of advanced laser matter-interaction in the context of warm dense matter research. WDM creation is
simulated using radiation-hydrodynamics to predict the thermodynamic state of the target during the exposure. We use a feedback loop
to optimize the optical laser pulse shape for low isentrope compression.


\printbibliography
\end{document}


